% ==================================================================================================================================
% Introduction

\minitoc  % Affiche la table des matières pour ce chapitre

Une fois définies les notions de boules, d'intérieur et d'adhérence, on peut maintenant "caractériser" 
des ensembles/parties en fonction des propriétés de leur adhérene/intérieur/frontière. 
Cela va nous permettre de définir les ouverts et les fermés, deux "catégories" d'ensembles essentielles pour 
la plupart des raisonnements analytiques de topologie. 

\vspace{0.3cm}

Ici, sauf mention contraire, nous nous plaçons dans un $\R$-espace vectoriel $E$. 
Pour toute partie $A$ de $E$, nous noterons son complémentaire dans $E$, $\overline{A}$, 
son adhérence $adh(A)$ et son intérieur $int(A)$. 

% ==================================================================================================================================
% Ouverts, Fermés

\section{Ouvert, Fermés}

\subsection{Définitions et Conventions}

\begin{definition}[Ensemble Ouvert]
    Soit $A \subseteq E$, on dit que $A$ est ouvert si $A = \int(A)$. 
    Autrement dit si pour tout élément de $A$, il existe une boule autour de cet élément entièrement 
    contenue dans $A$. 
\end{definition}

\begin{definition}[Ensemble Fermé]
    Soit $A \subseteq E$, on dit que $A$ est ouvert si $adh(A) = A$. 
\end{definition}

\begin{remark}
    Quelques conventions sur les ouverts et les fermés. 
    \begin{itemize}
        \item $ \forall a \in E, \forall r > 0 \quad B(a,r)$ est un ouvert et $\overline{B}(a,r)$ est un fermé. 
        \item $\emptyset$ et $\R^n$ dans $\R^n$ sont à la fois ouverts et fermés. 
        \item $[a,b[ \subset \R$ n'est ni ouvert, ni fermé.  
    \end{itemize}
\end{remark}

\subsection{Propriétés}

\begin{proposition}
    Soit $A \subset \R^n$, $A$ est ouvert ssi son complémentaire dans $\R^n$ est fermé. 
\end{proposition}

\begin{prop}[Réunion et Intersection]
    \begin{itemize}
        \item Une réunion quelconque d'ouverts est ouverte.
        \item Une intersection finie d'ouverts est ouverte.
        \item Une réunion finie de fermés est fermée.
        \item Une intersection quelconque de fermés est fermés
    \end{itemize}
\end{prop}

\begin{remark}[Moyen Mnémotechnique]
    Pour aider à la mémorisation, on peut s'aider de ces phrases :
    \begin{itemize}
        \item Les ouverts aiment s’étaler (réunion infinie), mais ils sont timides à se croiser (intersection finie).
        \item Les fermés aiment se serrer (intersection infinie), mais ne se dispersent pas trop (réunion finie).
    \end{itemize}
\end{remark}

\begin{example}[Réunion et Intersection]
    Quelques exemples pour retenir les propriétés :
    \begin{itemize}
        \item \textbf{Réunion infinie d'ouverts : } Soit $(A_n)_{n \in \N} = \left. \right] -\frac{1}{n}; \frac{1}{n} \left[ \right. $ une suite d'intervalles ouverts. 
        Alors la réunion infinie de tout ces intervalles reste ouverte :
            \[ \bigcup_{n = 1}^\infty  \left. \right] -\frac{1}{n}; \frac{1}{n} \left[ \right.  \quad \text{ouvert}\] 
        \item \textbf{Intersection Infinie de fermés : } De même, soit $(B_n)_{n \in \N} =  \left[-\frac{1}{n}; \frac{1}{n} \right]$
        une suite d'intervalles fermés. Alors leur intersection infinie reste fermée :
            \[ \bigcap_{n = 1}^\infty  \left[-\frac{1}{n}; \frac{1}{n} \right] \quad \text{fermé}\]
    \end{itemize}
\end{example}

\begin{proposition}
    Soit $A \subset E$, on dit que l'intérieur de $A$ est le plus grand ouvert contenu dans $A$ et 
    l'adhérence de $A$ est le plus petit fermé contenant $A$. 
\end{proposition}

\subsection{Ouverts/Fermés relativement}

\begin{definition}[Ouvert/Fermé relativement]
    Soient $A \subset E$ et $B \subset A$. On dit que $B$ est \textbf{ouvert relativement} à $A$ si il existe 
    une ouvert $V$ de $E$ tel que $B = A \cap V$. 

    D'autre part, on dit que $B$ est \textbf{fermé relativement} à $A$ si il existe un fermé 
    $U$ de $E$ tel que $B = A \cap U$. 
\end{definition}

\begin{comment}
\begin{remark}
    Illustrons ce nouveau concept d'ouvert/de fermé relativement à une partie. 
    Soit $A \subset E$ et $B,C \subset A$ respectivement un ouvert relativement à A et un fermé relativement à A. 

    \begin{center}
        \begin{tikzpicture}[scale=1.5]

            % Partie A (Cercle)
            \draw[name path=A] (0,0) circle (1.5) node[below right=1.6cm] {$A$};
            
            % Ouvert V (Ellipse en pointillé)
            \draw[dashed,name path=V] (1.25,0) ellipse (1.5cm and 0.25cm) node[below right=0.5cm] {$V$};
            
            % Fermé U (Ellipse pleine)
            \draw[name path=U] (-1.5,-1) ellipse (1cm and 0.75cm) node[below left=1cm] {$U$};
        
            % Trouver les intersections entre A et V, et remplir l'intersection
            \path[name intersections={of=A and V,by={[A-V]}}]
            \foreach \x in {[A-V]}{
                \fill[gray, opacity=0.3] (\x) circle (2pt);
            };
            
            % Trouver les intersections entre A et U, et remplir l'intersection
            \path[name intersections={of=A and U,by={[A-U]}}]
            \foreach \x in {[A-U]}{
                \fill[gray, opacity=0.5] (\x) circle (2pt);
            };
        
        \end{tikzpicture}
    \end{center}
\end{remark}
\end{comment}



% ==================================================================================================================================
% Ouverts, Fermés et fonctions continues 

\section{Ouvert, Fermés et Fonctions Continues}

Grâce à la notion d'ouverts et de fermés, on peut donner une caractérisation de la notion de continuité d'une fonction. 

\begin{definition}[Continuité, Rappel]
    Soit $E$ et $F$ deux $\R$-espaces vectoriel et $f : A \subset E \longrightarrow F$. 
    Soit $a \in A$ on dit que $f$ est continue en $a$ si 
        \[ \forall \varepsilon > 0, \exists \alpha > 0, \quad x \in A \text{ et } \| x - a \| \leqslant \alpha \Longrightarrow \| f(x) - f(a) \| \leqslant \varepsilon \] 
    Autrement dit, en terme de boules :
        \[ \forall \varepsilon > 0, \forall \alpha > 0, \quad f(A \cap B(a,\alpha)) \subset B(f(a), \varepsilon) \] 
\end{definition}

On peut alors énoncer une caractérisation de la continuité d'une fonction. 

\begin{proposition}[Continuité et Ouvert/Fermés]
    Soit $f : A \subset E \longrightarrow F$, $f$ est continue sur $A$ ssi pour tout ouvert $V$ de F, $f^{-1}(V)$ est ouvert 
    relativement à $A$. 
    On peut aussi dire que $f$ est continue sur $A$ ssi pour tout fermé $U$ de F, $f^{-1}(U)$ est fermé 
    relativement à $A$. 
\end{proposition}

\begin{remark}[Illustration]
    
\end{remark}



